\clearpage
\VerdierPhysicalPageAnchor{0051}%
\VerdierSemanticAnchor{ch1-1-5}
\subsection{\texorpdfstring{$\mathbb Z^n$}{Z n}-圏}

\VerdierSemanticAnchor{ch1-1-5-1}
\noindent\textbf{1.5.1.}\quad
$(\mathcal C,\mathbb Z^n,T,c)$ を $\mathbb Z^n$-圏、
$(\mathcal C',\mathbb Z^n,T',c')$ を加法的な $\mathbb Z^n$-圏とし、
\[
  (F,p)\colon(\mathcal C,\mathbb Z^n,T,c)
  \longrightarrow(\mathcal C',\mathbb Z^n,T',c')
\]
を正規化された余サイクル
$\alpha\in Z^2(\mathbb Z^n,\mathbb Z/2\mathbb Z)$ で捩られた
$\mathbb Z^n$-関手とする。$e_1,\ldots,e_n$ を $\mathbb Z^n$ の標準基底と
する。このとき、各 $i\in[1,n]$ に対して関手の同型

\VerdierSemanticAnchor{ch1-1-5-1-eq1}
\begin{equation}
  p(e_i)\colon FT(e_i)\longrightarrow T'(e_i)F
  \tag{1.5.1.1}
\end{equation}
があり、また $i\ne j$ である各対 $(i,j)$ に対して、次の図式は
$\alpha(e_i,e_j)\alpha^{-1}(e_j,e_i)$-可換である：

\VerdierSemanticAnchor{ch1-1-5-1-eq2}
\begin{center}
  \resizebox{0.98\textwidth}{!}{%
    $\begin{tikzcd}[ampersand replacement=\&,column sep=4.2em,row sep=3.3em]
      FT(e_i+e_j)
        \arrow[r,"{F*c(e_i,e_j)}"]
        \arrow[d,"{F*c(e_j,e_i)}"]
        \&
      FT(e_j)T(e_i)
        \arrow[r,"{p(e_j)*T(e_i)}"]
        \&
      T'(e_j)FT(e_i)
        \arrow[d,"{T'(e_j)*p(e_i)}"'] \\
      FT(e_i)T(e_j)
        \arrow[d,"{p(e_i)*T(e_j)}"]
        \& {} \&
      T'(e_j)T'(e_i)F
        \arrow[d,"{c'(e_i,e_j)^{-1}*F}"'] \\
      T'(e_i)FT(e_j)
        \arrow[r,"{T'(e_i)*p(e_j)}"']
        \&
      T'(e_i)T'(e_j)F
        \arrow[r,"{c'(e_j,e_i)^{-1}*F}"']
        \&
      T'(e_i+e_j)F
    \end{tikzcd}$%
  }\;.
  \hfill\textup{(1.5.1.2)}
\end{center}

\VerdierSemanticAnchor{ch1-1-5-2}
\noindent\textbf{命題 1.5.2.}\quad
\emph{$F\colon\mathcal C\to\mathcal C'$ を関手とし、各 $i\in[1,n]$ に
対して関手の同型
\[
  p_i\colon FT(e_i)\longrightarrow T'(e_i)F
\]
が与えられているとする。各対 $i\ne j$ について、図式 (1.5.1.2) の
$p(e_i)$ を $p_i$ に置き換えた図式が
$\alpha(e_i,e_j)\alpha^{-1}(e_j,e_i)$-可換であるならば、余サイクル
$\alpha$ で捩られた $\mathbb Z^n$-関手
\[
  (F,p)\colon(\mathcal C,\mathbb Z^n,T,c)
  \longrightarrow(\mathcal C',\mathbb Z^n,T',c')
\]
で
\[
  p(e_i)=p_i,
  \qquad i\in[1,n]
\]
を満たすものがただ一つ存在する。}

\VerdierSemanticAnchor{ch1-1-5-2-proof-intro}
一意性は図式 (1.4.3.11) を調べれば証明できる。この確認は読者に委ねる。
ここでは存在証明の指針だけを与える。証明は整数 $n$ に関する帰納法で行う。
